\documentclass[a4paper,12pt]{article}
\usepackage[utf8]{inputenc}
\usepackage[T2A]{fontenc}
\usepackage[russian]{babel}
\usepackage{amsmath, amssymb, amsthm}
\usepackage{geometry}
\usepackage{booktabs}
\usepackage{array}
\usepackage{hyperref}
\usepackage{xcolor}
\usepackage{graphicx}
\usepackage{float}
\usepackage{caption}
\usepackage{subcaption}
\usepackage{tikz}
\usetikzlibrary{arrows.meta, positioning}

\geometry{left=2.5cm, right=2cm, top=2cm, bottom=2cm}

\hypersetup{
    colorlinks=true,
    linkcolor=blue!70!black,
    urlcolor=blue!70!black,
    citecolor=red!70!black
}

\title{Лабораторная работа №5\\Задача регрессии}

\author{
    Дмитрий Коваль, Георгий Лаврентьев\\
    группа М3232
}
\date{2 июня 2026 г.}

\newcommand{\R}{\mathbb{R}}
\newcommand{\norm}[1]{\left\|#1\right\|}
\newcommand{\grad}{\nabla}
\newcommand{\eps}{\varepsilon}
\newcommand{\inner}[2]{\left\langle #1, #2 \right\rangle}

\theoremstyle{definition}
\newtheorem{definition}{Определение}[section]
\newtheorem{theorem}{Теорема}[section]
\newtheorem{remark}{Замечание}[section]

\begin{document}
\maketitle
\vspace{-20pt}

% ═════════════════════════════════════════════════════════════
\section{Постановка задачи}
% ═════════════════════════════════════════════════════════════

Рассматривается задача восстановления регрессионной зависимости по зашумлённым данным. По выборке
$\{(x_i, y_i)\}_{i=1}^{m}$ строится параметрическая модель $\widehat{y}(x, w)$, а вектор параметров
$w \in \R^p$ выбирается из условия минимизации эмпирического риска
\[
    Q(w) = \frac{1}{m} \sum_{i=1}^{m} \bigl(\widehat{y}(x_i, w) - y_i\bigr)^2 \longrightarrow \min_{w}.
\]
При наличии регуляризации минимизируется полная функция потерь
\[
    L(w) = Q(w) + \lambda_1 \norm{w}_1 + \lambda_2 \norm{w}_2^2.
\]

В работе реализованы и исследованы две модели регрессии (линейная и полиномиальная),
четыре варианта функции потерь (без регуляризации, $L_1$, $L_2$ и Elastic Net) и пять методов
оптимизации:
\begin{enumerate}
    \item аналитическое решение линейной регрессии через оценки средних (одномерный случай);
    \item стохастический градиентный спуск (SGD);
    \item mini-batch градиентный спуск;
    \item метод Гаусса--Ньютона;
    \item метод Гаусса--Ньютона с регуляризацией Левенберга--Марквардта.
\end{enumerate}

\paragraph{Единый интерфейс и подсчёт затрат.}
Модели имеют единый интерфейс (\texttt{design} — построение матрицы плана, \texttt{predict} —
предсказание), а оптимизаторы — единую сигнатуру запуска вида
\[
    \texttt{optimizer}(\texttt{problem},\ w_0,\ \dots) \longrightarrow \texttt{RegressionResult},
\]
где \texttt{problem} инкапсулирует матрицу плана $\Phi$, целевой вектор $y$ и параметры
регуляризации, а \texttt{RegressionResult} содержит найденный вектор $w$, итоговые значения
$L$, $Q$ и регуляризатора, а также историю этих величин и учётные показатели:
\begin{enumerate}
    \item число итераций / эпох;
    \item число вычислений градиента (число обработанных батчей);
    \item время работы;
    \item причину остановки.
\end{enumerate}
Для SGD и mini-batch дополнительно фиксируются размер батча, число эпох и правило выбора шага.
Подсчёт вызовов оракулов доступен через декоратор-счётчик \texttt{CallCounter}; все
векторно-матричные операции выполнены средствами \texttt{NumPy}.

\paragraph{Причины остановки.}
Методы возвращают различимые статусы завершения:
\begin{itemize}
    \item \textit{Сходимость по градиенту} — выполнено $\norm{\grad L(w_k)} \leq \eps$;
    \item \textit{Сходимость по функции потерь} — относительное изменение потерь за эпоху
          опустилось ниже порога, $|L_{k+1} - L_k| \leq \texttt{tol}\,(1 + |L_k|)$;
    \item \textit{Сходимость (шаг не уменьшает потери)} — в методе Левенберга--Марквардта, например, параметр демпфирования вырос настолько, что ни один шаг не уменьшает $L$, при малом градиенте;
    \item \textit{Достигнут лимит эпох / итераций} — исчерпан бюджет шагов;
    \item \textit{Расходимость} — значение потерь перестало быть конечным;
    \item \textit{Аналитическое решение} — решение получено в замкнутой форме.
\end{itemize}

% ═════════════════════════════════════════════════════════════
\section{Генерация данных}
% ═════════════════════════════════════════════════════════════

Данные генерируются на основе аналитических функций с добавлением нормального шума
\[
    y_i = f_{\text{true}}(x_i) + \xi_i, \qquad \xi_i \sim \mathcal{N}(0, \sigma^2),
\]
где точки $x_i$ равномерно распределены на отрезке $[-3, 3]$, а объём каждой выборки $m = 160$
($\geq 100$ согласно условию). Рассматриваются два набора данных.

\begin{description}
    \item[Почти линейная зависимость.]
        \[
            f_{\text{true}}(x) = 0.8\,x - 0.5 + 0.1 \sin(3x), \qquad \sigma = 0.35.
        \]
        Слагаемое $0.1\sin(3x)$ задаёт слабое отклонение от прямой.
    \item[Сильно нелинейная зависимость.]
        \[
            f_{\text{true}}(x) = \sin(2x) + 0.15\,x^3 - 0.5\,x
            + 1.5\,\exp\!\Bigl(-\tfrac{(x-1)^2}{0.1}\Bigr), \qquad \sigma = 0.45.
        \]
        Здесь сочетаются периодическая часть $\sin(2x)$, полиномиальная часть $0.15x^3 - 0.5x$
        и локальный пик в окрестности $x = 1$.
\end{description}
% Истинная зависимость, дисперсия шума и диапазон $x$ для обеих задач явно зафиксированы выше.

% ═════════════════════════════════════════════════════════════
\section{Модели регрессии}
% ═════════════════════════════════════════════════════════════

\begin{definition}[Полиномиальная модель]
    Моделью полиномиальной регрессии степени $d$ называется
    \[
        \widehat{y}(x, w) = \sum_{j=0}^{d} w_j\, \varphi_j(x),
    \]
    где $\varphi_j$ — признаки, построенные из $x$. Линейная регрессия соответствует $d = 1$.
\end{definition}

Модель линейна по параметрам $w$, поэтому на выборке предсказания записываются через
\emph{матрицу плана} $\Phi \in \R^{m \times (d+1)}$ со строками
$\varphi(x_i) = (1, z_i, z_i^2, \dots, z_i^d)$ и
\[
    \widehat{y} = \Phi w.
\]
Согласно рекомендации условия, входной признак предварительно нормируется,
$z = (x - \bar{x}) / s_x$, после чего столбцы матрицы плана дополнительно центрируются и
масштабируются к единичной дисперсии (кроме свободного столбца). Это существенно уменьшает
число обусловленности задачи: для нормированного базиса при степени $d=5$ оно составляет
$\kappa \approx 5\cdot 10^2$, тогда как для «сырого» базиса Вандермонда той же степени
$\kappa$ на несколько порядков больше (см. раздел~\ref{sec:methods-compare}).

% ═════════════════════════════════════════════════════════════
\section{Функции потерь и регуляризация}
% ═════════════════════════════════════════════════════════════

% Эмпирический риск и его градиент для модели, линейной по параметрам, имеют вид
% \[
%     Q(w) = \frac{1}{m} \norm{\Phi w - y}_2^2, \qquad
%     \grad Q(w) = \frac{2}{m}\, \Phi^\top (\Phi w - y).
% \]
Регуляризованная функция потерь
\[
    L(w) = Q(w) + \lambda_1 \sum_{j \geq 1} |w_j| + \lambda_2 \sum_{j \geq 1} w_j^2
\]
% охватывает все четыре требуемых варианта: $\lambda_1 = \lambda_2 = 0$ (без регуляризации),
% только $\lambda_1 > 0$ ($L_1$), только $\lambda_2 > 0$ ($L_2$) и $\lambda_1, \lambda_2 > 0$
% (Elastic Net). Свободный коэффициент $w_0$ не штрафуется.

\begin{remark}[Сглаживание $L_1$]
    Модуль недифференцируем в нуле, поэтому, как допускает условие, используется гладкая
    аппроксимация
    \[
        |w_j| \approx \sqrt{w_j^2 + \eps^2}, \qquad
        \frac{\partial}{\partial w_j}\sqrt{w_j^2 + \eps^2} = \frac{w_j}{\sqrt{w_j^2 + \eps^2}},
    \]
    с малым $\eps = 10^{-6}$. Это позволяет применять к $L_1$ и Elastic Net те же градиентные
    методы, что и к гладким потерям.
\end{remark}

% ═════════════════════════════════════════════════════════════
\section{Методы оптимизации}\label{sec:methods}
% ═════════════════════════════════════════════════════════════

\subsection{Аналитическое решение (одномерный случай)}

Для линейной регрессии с одним признаком оптимальные коэффициенты выражаются через оценки средних:
\[
    \widehat{w}_1 = \frac{\sum_{i} (z_i - \bar{z})(y_i - \bar{y})}{\sum_{i} (z_i - \bar{z})^2},
    \qquad
    \widehat{w}_0 = \bar{y} - \widehat{w}_1 \bar{z}.
\]
Это точное решение задачи наименьших квадратов в одномерном случае, не требующее итераций.

\subsection{Градиентный спуск, SGD и mini-batch}

Итерационные методы записываются в форме
\[
    w_{k+1} = w_k - \alpha_k\, p_k,
\]
где для базового спуска $p_k = \grad L(w_k)$. В стохастической постановке полный градиент
заменяется его оценкой по случайному подмножеству (батчу) $B \subseteq \{1, \dots, m\}$:
\[
    \widehat{\grad} Q(w) = \frac{2}{|B|}\, \Phi_B^\top (\Phi_B w - y_B),
\]
к которой добавляется точный градиент регуляризатора. Метод SGD соответствует $|B| = 1$,
mini-batch — промежуточным значениям $1 < |B| < m$, а полный батч $|B| = m$ даёт обычный
градиентный спуск. Одна \emph{эпоха} — это один проход по случайно перемешанной выборке;
число вычислений градиента равно числу обработанных батчей.

\paragraph{Правило выбора шага.}
Используется убывающий по времени шаг
\[
    \alpha_t = \frac{\alpha_0}{1 + \gamma t},
\]
где $t$ — номер обновления, $\alpha_0$ — начальный шаг, $\gamma$ — скорость затухания
(поддерживается и постоянный шаг). Затухание обеспечивает устойчивость стохастических методов
вблизи минимума, где шум оценки градиента не позволяет точно остановиться.

\subsection{Метод Гаусса--Ньютона}

Для задачи наименьших квадратов с вектором $r(w) = \Phi w - y$ и якобианом
$J = \partial r / \partial w = \Phi$ (постоянным, так как модель линейна по параметрам) метод
Гаусса--Ньютона приближает гессиан риска матрицей $\tfrac{2}{m} J^\top J$ и делает шаг
\[
    w_{k+1} = w_k - H^{-1} \grad L(w_k), \qquad
    H = \frac{2}{m}\, \Phi^\top \Phi + H_{\text{рег}},
\]
где $H_{\text{рег}}$ — диагональный вклад гладкого регуляризатора. Для модели, линейной по
параметрам, шаг Гаусса--Ньютона совпадает с решением нормальных уравнений, поэтому метод
сходится за одну итерацию (вторая итерация лишь подтверждает $\norm{\grad L} \leq \eps$).

\subsection{Метод Левенберга--Марквардта}

Левенберг--Марквардт демпфирует систему Гаусса--Ньютона параметром $\mu \geq 0$:
\[
    (H + \mu I)\, \Delta = \grad L(w_k), \qquad w_{k+1} = w_k - \Delta.
\]
Согласно условию, реализовано правило адаптации $\mu$: если пробный шаг уменьшает потери, он
принимается, а $\mu$ уменьшается ($\mu \leftarrow \mu / \nu$); если нет — шаг отклоняется, а
$\mu$ увеличивается ($\mu \leftarrow \mu \cdot \nu$) до тех пор, пока шаг не станет успешным
(в работе $\mu_0 = 10^{-2}$, $\nu = 3$). При $\mu \to 0$ метод приближается к Гауссу--Ньютону,
при больших $\mu$ — к мелкому шагу вдоль антиградиента. Демпфирование делает метод устойчивым
на плохо обусловленных задачах, где прямое решение нормальных уравнений ненадёжно.

% ═════════════════════════════════════════════════════════════
\section{Экспериментальная часть}
% ═════════════════════════════════════════════════════════════

\subsection{Задание 1: базовые модели без регуляризации}

Для обеих задач обучены линейная и полиномиальные (степеней $2$--$5$) модели всеми применимыми
методами. В таблицах 1 и 2 приведён итоговый
эмпирический риск $Q$.

\input{results/tables/base_linear_risk.tex}
\input{results/tables/base_nonlinear_risk.tex}

Аналитическое решение, Гаусс--Ньютон и Левенберг--Марквардт достигают одного и того же минимума
наименьших квадратов (значения совпадают до четырёх знаков), что подтверждает корректность
реализаций. Для почти линейной задачи рост степени почти не уменьшает риск: зависимость и так
близка к прямой. Для нелинейной задачи риск резко падает при переходе к степени $5$
($0.73 \to 0.29$): только полином достаточной степени способен описать локальный пик.

Метод mini-batch практически достигает оптимума наименьших квадратов, тогда как SGD при том же
бюджете эпох заметно отстаёт на высокой степени — это закономерное следствие плохой
обусловленности (см.\ рис.~\ref{fig:nonlinear_methods}, зелёная кривая).

\begin{figure}[H]
    \centering
    \begin{subfigure}{0.49\textwidth}
        \includegraphics[width=\textwidth]{results/plots/fig_linear_degrees.png}
        \caption{Почти линейная: полиномы разных степеней.}
    \end{subfigure}
    \hfill
    \begin{subfigure}{0.49\textwidth}
        \includegraphics[width=\textwidth]{results/plots/fig_nonlinear_degrees.png}
        \caption{Нелинейная: полиномы разных степеней.}
    \end{subfigure}
    \caption{Аппроксимация полиномами разных степеней (решение Гаусса--Ньютона). Для нелинейной
             задачи только степень $5$ улавливает форму зависимости.}
    \label{fig:degrees}
\end{figure}

\begin{figure}[H]
    \centering
    \begin{subfigure}{0.49\textwidth}
        \includegraphics[width=\textwidth]{results/plots/fig_linear_methods.png}
        \caption{Почти линейная, степень $1$.}
    \end{subfigure}
    \hfill
    \begin{subfigure}{0.49\textwidth}
        \includegraphics[width=\textwidth]{results/plots/fig_nonlinear_methods.png}
        \caption{Нелинейная, степень $5$.}
    \end{subfigure}
    \caption{Сравнение методов на одной модели. Точные методы (аналит., GN, LM) и mini-batch
             совпадают; SGD на нелинейной задаче недообучен при том же числе эпох.}
    \label{fig:nonlinear_methods}
\end{figure}

\begin{figure}[H]
    \centering
    \begin{subfigure}{0.49\textwidth}
        \includegraphics[width=\textwidth]{results/plots/fig_linear_loss.png}
        \caption{Почти линейная.}
    \end{subfigure}
    \hfill
    \begin{subfigure}{0.49\textwidth}
        \includegraphics[width=\textwidth]{results/plots/fig_nonlinear_loss.png}
        \caption{Нелинейная.}
    \end{subfigure}
    \caption{Динамика функции потерь по эпохам для итерационных методов SGD и mini-batch.}
    \label{fig:baseline_loss}
\end{figure}

\subsection{Задание 2: влияние размера batch}

На нелинейной задаче с полиномом степени $5$ метод mini-batch запущен при размерах батча
$1, 4, 8, 16, 32, 64, m$ с одинаковыми шагом, затуханием и числом эпох. Итоговые показатели
сведены в таблицу~\ref{tab:batch_size}.

\input{results/tables/batch_size.tex}

\begin{figure}[H]
    \centering
    \begin{subfigure}{0.49\textwidth}
        \includegraphics[width=\textwidth]{results/plots/fig_batch_epoch.png}
        \caption{По эпохам.}
    \end{subfigure}
    \hfill
    \begin{subfigure}{0.49\textwidth}
        \includegraphics[width=\textwidth]{results/plots/fig_batch_grad.png}
        \caption{По числу вычислений градиента.}
    \end{subfigure}
    \caption{Динамика потерь при разных размерах батча.}
    \label{fig:batch}
\end{figure}

При фиксированном числе эпох меньший батч даёт больше обновлений за эпоху и потому быстрее
снижает потери \emph{по эпохам}: за $200$ эпох $|B|=1$ доходит до риска $\approx 0.43$, а полный
батч — лишь до $\approx 0.64$. Платой за это служат шум траектории и существенно большее число
вычислений градиента ($32000$ против $200$). В пересчёте \emph{на одну эпоху} крупные батчи
дешевле и стабильнее. Разумным компромиссом для данной задачи оказались средние размеры
$|B| = 8$--$16$: они сочетают приемлемую скорость убывания потерь с устойчивостью траектории.

\subsection{Задание 3: влияние регуляризации}

Переобучение демонстрируется на нелинейной задаче с полиномом высокой степени ($d = 9$).
Базовая модель без регуляризации сравнивается с $L_1$, $L_2$ и Elastic Net при нескольких
значениях параметров (таблица 4); модули коэффициентов до и после
регуляризации приведены в таблице 5.

\input{results/tables/reg_summary.tex}
\input{results/tables/reg_coeffs.tex}

\begin{figure}[H]
    \centering
    \begin{subfigure}{0.32\textwidth}
        \includegraphics[width=\textwidth]{results/plots/fig_reg_dyn_l1.png}
        \caption{$L_1$.}
    \end{subfigure}
    \hfill
    \begin{subfigure}{0.32\textwidth}
        \includegraphics[width=\textwidth]{results/plots/fig_reg_dyn_l2.png}
        \caption{$L_2$.}
    \end{subfigure}
    \hfill
    \begin{subfigure}{0.32\textwidth}
        \includegraphics[width=\textwidth]{results/plots/fig_reg_dyn_en.png}
        \caption{Elastic Net.}
    \end{subfigure}
    \caption{Динамика полных потерь $L$, эмпирического риска $Q$ и регуляризатора по эпохам.}
    \label{fig:reg_dyn}
\end{figure}

\begin{figure}[H]
    \centering
    \begin{subfigure}{0.49\textwidth}
        \includegraphics[width=\textwidth]{results/plots/fig_reg_fit.png}
        \caption{Аппроксимация.}
    \end{subfigure}
    \hfill
    \begin{subfigure}{0.49\textwidth}
        \includegraphics[width=\textwidth]{results/plots/fig_reg_coeffs.png}
        \caption{Модули коэффициентов (лог. шкала).}
    \end{subfigure}
    \caption{Влияние регуляризации на аппроксимацию и на структуру коэффициентов модели степени $9$.}
    \label{fig:reg}
\end{figure}

Без регуляризации модель степени $9$ переобучается: коэффициенты достигают величин порядка
$10$--$25$, а кривая повторяет шум (рис.~\ref{fig:reg}). Все три типа регуляризации резко
уменьшают модули коэффициентов и сглаживают аппроксимацию. Различие между типами проявляется в
структуре решения: $L_1$ и Elastic Net обнуляют часть коэффициентов (в таблице~\ref{tab:reg_coeffs}
у них $5$--$6$ ненулевых весов из $9$), давая разрежённое решение, тогда как $L_2$ равномерно
сжимает все коэффициенты, не зануляя их. Это согласуется с известными свойствами $L_1$- и
$L_2$-штрафов.

\subsection{Задание 4: сравнение методов оптимизации}\label{sec:methods-compare}

Сводное сравнение методов по итоговому риску, числу итераций/эпох, числу вычислений градиента и
времени приведено в таблице 6.

\input{results/tables/methods_compare.tex}

Точные методы (аналит., GN, LM) достигают минимума за единицы итераций и доли миллисекунды.
Первопорядковые методы требуют сотен эпох и тысяч--десятков тысяч вычислений градиента; SGD при
этом наименее эффективен на плохо обусловленной задаче. Статус «$\times$» у SGD и mini-batch
означает исчерпание бюджета эпох (выход на шумовой уровень), а не расходимость: достигнутый риск
близок к оптимуму.

\paragraph{Гаусс--Ньютон против Левенберга--Марквардта.}
Чтобы выявить роль обусловленности, методы сравнены на «сыром» (ненормированном по столбцам)
базисе Вандермонда степеней $8$ и $12$, где число обусловленности гессиана достигает
$\kappa \sim 10^{5}$ и $10^{9}$ (таблица 7).

\input{results/tables/gn_lm_illcond.tex}

\begin{figure}[H]
    \centering
    \includegraphics[width=0.62\textwidth]{results/plots/fig_gn_lm.png}
    \caption{Сходимость Гаусса--Ньютона и Левенберга--Марквардта на необусловленном базисе.}
    \label{fig:gn_lm}
\end{figure}

Оба метода достигают близкого итогового риска, однако ведут себя по-разному. Гаусс--Ньютон
делает один шаг прямого решения системы: при $\kappa \sim 10^{9}$ этот шаг численно ненадёжен и
полностью полагается на устойчивость библиотечного решателя. Левенберг--Марквардт за счёт
демпфирования $\mu$ строит последовательность гарантированно убывающих по $L$ шагов и потому
устойчив к плохой обусловленности; платой служит большее число итераций ($9$--$12$ против $2$).

% ═════════════════════════════════════════════════════════════
\section{Выводы}
% ═════════════════════════════════════════════════════════════

\begin{enumerate}
    \item Реализованы все требуемые компоненты: генератор двух наборов данных, линейная и
          полиномиальная модели, четыре варианта функции потерь ($L_1$, $L_2$, Elastic Net и без
          регуляризации) и пять методов оптимизации с единым интерфейсом и сохранением истории
          потерь, риска и регуляризатора.
    \item Аналитическое решение, Гаусс--Ньютон и Левенберг--Марквардт дают идентичный минимум
          наименьших квадратов; для модели, линейной по параметрам, Гаусс--Ньютон эквивалентен
          нормальным уравнениям и сходится за одну итерацию.
    \item Сложность модели нужно согласовывать с данными: нелинейную зависимость с локальным
          пиком описывает лишь полином степени $\geq 5$, а избыточная степень $9$ без
          регуляризации приводит к переобучению.
    \item Размер батча управляет компромиссом «скорость--устойчивость--стоимость»: малые батчи
          быстрее убывают по эпохам, но шумны и дороги по числу вычислений градиента; наилучшими
          для рассмотренной задачи оказались средние размеры $|B| = 8$--$16$.
    \item Регуляризация подавляет переобучение и уменьшает модули коэффициентов; $L_1$ и Elastic
          Net дополнительно разрежают решение, тогда как $L_2$ равномерно сжимает веса.
    \item Обусловленность — ключевой фактор: при росте степени полинома $\kappa$ быстро растёт,
          из-за чего первопорядковые методы замедляются, прямой шаг Гаусса--Ньютона становится
          ненадёжным, а демпфирование Левенберга--Марквардта обеспечивает устойчивую сходимость.
\end{enumerate}

\end{document}
